\section{Présentation d'Enedis}
\subsection{Missions et activités principales}
% Description du rôle d'Enedis en tant que gestionnaire du réseau de distribution d'électricité
% Chiffres clés : nombre d'employés, couverture territoriale, etc.
\lipsum[1]

\subsection{Organisation générale}
% Structure organisationnelle d'Enedis
% Les différentes directions et leur rôle
% Place des Directions Régionales (DR) dans l'organisation
\lipsum[2]

\section{Le Centre De Service : la cellule d'innovation informatique}
\subsection{Rôle et missions}
% Présentation du CDS et son positionnement dans l'entreprise
% Missions spécifiques : innovation, développement d'applications, etc.
\lipsum[3]

\subsection{Organisation de l'équipe}
% Structure de l'équipe : composition, rôles
% Méthodologies de travail
% Culture d'entreprise
\lipsum[4]

\section{Environnement technique}
\subsection{Écosystème technologique}
% Technologies principales utilisées
% Stack technique : NestJS, Angular, Prisma, etc.
% Outils de développement et de collaboration
\lipsum[5]

\subsection{Contraintes et spécificités}
% Contraintes liées au secteur de l'énergie
% Enjeux de sécurité et de confidentialité
% Intégration avec les systèmes existants d'Enedis
\lipsum[6]

\section{Positionnement du projet dans l'organisation}
\subsection{Contexte et enjeux RH}
% Enjeux RH spécifiques à Enedis
% Importance du projet dans la stratégie RH
\lipsum[7]

\subsection{Parties prenantes}
% Identification des parties prenantes principales
% Rôles et attentes des différents acteurs
\lipsum[8]